\section{Conclusion}
\label{sec:conclusion}

This thesis built and validated the terminal stage of underwater docking using passive markers and classical computer vision, the approach available to low-cost platforms. The system consists of multi-marker ArUco fusion, an error-state Kalman filter with two dock-motion models, and a coarse-to-fine docking state machine. The full perception pipeline runs in the loop, and the vehicle never receives ground truth.

In simulation, the system docks onto a moving target. The baseline configuration docked at every sway period down to 8~s at 0.1~m amplitude, which covers moored docks at typical continental-shelf depths in the mean New South Wales sea state. The velocity feedforward fails at short periods because the vehicle's own motion leaks into the dock measurement and is integrated into a phantom dock velocity. Tuning cannot fix this failure, the system's own health signals cannot see it, and the baseline is immune because it never estimates velocity.

On real underwater imagery, detection needs about 27 pixels of apparent marker size, and the maximum detection range is about 35 times the marker side. Detection collapses once the total optical depth exceeds about 2. In the measured pool water, these numbers set the marker for 5~m acquisition at 182~mm (154 to 207~mm).

Together, the two parts bound where this class of system works, from the motion side and from the sensing side. They also supply the numbers a dock designer needs, namely a capture funnel of at least 0.15~m radius tolerating about 0.1~m/s closing speed, and the marker sizing rule above. Both quantities come from measurement. In addition, no stage of the pipeline depends on a trained model, so the measured limits are properties of the approach itself.

\subsection{Future work}
\label{sec:future-work}

The first priority is to fix the velocity estimator. It breaks because the dock measurement is built through the vehicle's own pose, so the vehicle's motion ends up inside the measurement. There are two ways to remove this. The filter could take in the vehicle odometry, so it knows which part of the motion belongs to the vehicle. Alternatively, the dock could be estimated directly relative to the vehicle, so the vehicle's world position is never involved. Either fix should be tested against the same sway sweep, and it still has to work with the 240~ms lag, since the instability appears once the lag becomes large compared to the sway period.

The natural next step is hardware testing with the combined work of Part~A and Part~B. Part~B measured the camera side in an open loop, while Part~A closed the loop only in simulation. A pool trial with the real BlueROV2, the printed marker board and the full docking stack would test both halves together for the first time. The 240~ms lag should be measured again on the real vehicle early on, since the failure period can be predicted from it before any docking is attempted.

On the measurement side, the camera should be calibrated underwater. This would remove the $\pm$10\% range uncertainty that dominates Part~B's error budget, since the metric ranges and the marker sizing rule would then rest on a measured focal length instead of the in-air factory value.

The detection measurements should also be repeated in genuinely turbid water. The synthetic turbidity adds contrast loss but no scattering blur, so the high-turbidity results are an upper bound. A repeat in a dosed and isolated tank or a naturally murky site would show whether the $\sim$27~px budget and the $\tau \approx 2$ cliff still hold with real scattering.

Finally, a sway-aware dock motion model may recover the feedforward benefit without the ego-motion problem. The dock moves in a nearly perfect sine at one wave period, so a model built for that motion is much more constrained than a general constant-velocity model. A more constrained model is harder to corrupt with motion that does not look like a sine, though this remains to be shown against the same sweep.
